\section{Related Work}
  
%Many existing studies have been discussed throughout the paper. In this section,
%we mainly focus on discussing other related work in the following two categories.

\vspace{-.1in}
\paragraph{Benchmarks.}
Many areas have well-established benchmarks for performance evaluation and comparison,
such as the TPC series~\cite{tpcurl}, SmallBank~\cite{alomari2008cost}, YCSB~\cite{Cooper2010YCSB},
fio~\cite{fio-repo}, the SPEC series~\cite{specurl},
HPL~\cite{hplurl}, HPCG~\cite{hpcg}, Graph500~\cite{graph500url},
MLPerf~\cite{mlperf-train,mlperf-infer}, DAWNBench~\cite{DAWNBench} and many others.
Most of these benchmarks have restrictions about what parameters can be tuned and what values can be used
for these parameters.
While this work focuses on
TPC-C and YCSB, we plan to perform a similar study on other
benchmarks in the future.

%\noindentparagraph{Database and key-value stores.}

\vspace{-.1in}
\paragraph{Reproducibility.} 
To improve code quality and facilitate future research, our community has encouraged
research papers to verify their reproducibility through artifact evaluation~\cite{osdi-repro,sosp-repro,eurosys-repro,vldb-repro,sigmod-repro}.
In the meantime, many works study how to improve the reproducibility of research
prototypes. For example, CloudLab studies the performance variability of its experiments
and proposes guidelines to reduce such variability~\cite{Maricq2018Variability};
a recent work studies how cloud-based big data workload can be affected by network variability~\cite{alex-nsdi20};
Fursin discussed the challenges and solutions from his experience of reproducing 150+ research
papers~\cite{fursin-2021};
Ursprung~\cite{lukas-vldb20} studies how to improve the reproducibility of data science workloads.
%The key idea of Ursprung is to transparently capture provenance and build lineage to automatically track static and runtime configuration parameters of data science pipelines. With Ursprung, the authors hope to improve the reproducibility of data science workloads.
%
Compared to these efforts, our work further argues that, for the purpose of
improving code quality and facilitating future research, we should encourage
extensive evaluation to cover a wider range of settings.

\vspace{-.1in}
\paragraph{Regression/Machine learning for systems.} 
%Regression analysis aims to estimate the relation between one dependent variable and multiple independent variables. 
%Regression is usually used to predict %or forecast 
%the dependent values with the independent values by a linear or non-linear regression model. 
Regression or machine learning based methods have been gaining momentum in systems designs and optimizations.
%The most widely used regression model is linear regression, which represents a line that fits the data. In addition to linear regression model, non-linear regression is more complex, which substantially overlaps with the field of machine learning. 
For instance, many studies~\cite{yu-dac-asplos18,mr-advisor,rfhoc} propose to use regression-based algorithms to tune systems configurations to achieve better performance. 
Recent systems research successfully exploits machine learning methods for resource management~\cite{Quasar,JouleGuard,LLAMA,CALOREE}, scheduling~\cite{Paragon,Decima}, I/O optimizations~\cite{Vivace,linnos}, and so on.
Compared to these efforts, our work demonstrates that regression can be a feasible solution to reduce the cost of extensive evaluation.

%works under a variety of settings, some of which are beyond the ones used in the original papers. As discussed in Section~\ref{sec:questions},
%to support such extensive comparison, we suggest to introduce the idea of ``supporting extensive testing'' into artifact evaluation.
%\vspace{-.1in}
%\paragraph{Extensive evaluation.} A number of works did extensive experiments to understand the impact of different concurrency control
%features~\cite{Harding2017EvalDistributed,Wu2017EvalMem,Lim2017Cicada}.
%This work studies the feasibility of extensive evaluation in general, how to visualize evaluation results,
%and the benefits of performing extensive evaluation.

%Since they also use TPC-C and YCSB, their observations overlap with ours to some extent (e.g., contention level matters).
%However, because they have a different goal compared to this paper, there are a few key differences: first, for an apple-to-apple comparison,
%they re-implement different concurrency control features under the same framework; this paper, instead, tries to understand whether
%changing experiment settings will change the conclusion of the original work and thus has to re-produce original works. Second, they focus more on
%tuning concurrency control features but less on tuning benchmark parameters and our work has the opposite focus. Finally, our
%paper covers additional systems which explore different designs (e.g. Janus, HERD, etc).
%Among the works we studied,
%Cicada~\cite{Lim2017Cicada} did an extensive comparison, but only on in-memory database engines.
%However, the major difference from our work is we are focusing on how to perform fair comparisons among systems by properly using popular benchmarks in our community. Note that, this paper also shows that the reproducibility of our studied systems are reasonably good.

%found that our community regularly disregards the variability when running experiments in the cloud. Thus, it
\vspace{-.1in}
\paragraph{Visualizing high-dimensional data.}
Visualizing high-dimensional data is a classic research topic with a large volume of work.
However, most of them focus on how to find the structure within high-dimensional data
and map it onto a 2D figure. This paper has a completely different goal and thus cannot
re-use these works directly. We have tried two popular algorithms in this area t-SNE~\cite{tSNE} and UMAP~\cite{UMAP} on our data 
and do not find a meaningful output. 
%t-SNE and UMAP both map high-dimensional data into a low dimension space, but with different techniques. 
%tSNE uses the t-distributed variant of Stochastic Neighbor Embedding~\cite{SNE}, 
%while UMAP builds upon mathematical foundations related to the work of Belkin and Niyogi on
%Laplacian eigenmaps~\cite{UMAPbackground}.
